%% =====================================================================
% Capítulo: Servidor de correo para el Escenario 1 (fuga por correo)
% Se incluye desde el main con \include{servidor} (sin extensión).
%
% Todo el diseño sale de styles/mistilo.sty (que ya carga listings,
% booktabs, caption, float, xcolor, tcolorbox, etc.). Este capítulo solo
% usa lo que ese estilo define:
%   - listados:  style=terminal
%   - avisos:    entorno aviso
%   - tablas:    booktabs + caption
% Idioma y codificación los decide el main; hace falta:
%   \usepackage[T1]{fontenc}   % para que < y > salgan bien en \texttt
% =====================================================================

% Código en línea (admite _, $, ^, {}, <, > ... sin escapar)
\providecommand{\cod}[1]{\texttt{\detokenize{#1}}}

\chapter{Configuración del servidor de correo electrónico para el escenario de fuga de información por correo}
\label{cap:servidor}

El objetivo de este capítulo es explicar la creación del sistema de correo preparado para el primero de los escenarios detallados en el enunciado; si bien no entraremos en detalles muy técnicos, en caso de querer profundizar en el código fuente de la aplicación así como la información técnica y una guía de uso e instalación, se puede revisar en el \href{https://github.com/ALVPI/IFOR_P2/tree/main/servidor}{repositorio del grupo H} para esta práctica.


Para poder realizar este trabajo vamos a empler el protocolo SMPT según se especifica en el \href{https://www.rfc-editor.org/info/rfc5321/}{RFC 5321} adaptándolo para nuestro entorno el cual van a ser dos VMs (máquinas virtuales) con el sistema operativo Windows 11.
\section{Funcionalidades}
Nuestro sevidor de correo y por consecuencia su portal web tiene van a realizar las siguientes funcionalidades:
\begin{enumerate}
  \item permita enviar \cod{Confidencial.pdf} como adjunto entre dos buzones
        del grupo \cod{pepe@empresa.local} $\rightarrow$
        \cod{ana@empresa.local}; los dos usuarios por defecto que hemos creado en el propio servidor.
  \item dejare \textbf{evidencia analizable} en disco (\cod{.eml} completo con
        cabecera \cod{Received}, adjunto extraído, log de auditoría);
\end{enumerate}
% ---------------------------------------------------------------------
\section{Decisiones técnicas y por qué}
Una vez tenemos claro cuál es el objetivo a cumplir, hemos tenido que tomar una serie de decisiones sobre el \textit{stack} tecnológico, es decir, protocolos, leguajes, librerias que vamos a emplear siguiendo lo especificado en el RFC correspondiente 
\begin{table}[H]
\centering
\small
\renewcommand{\arraystretch}{1.3}
\begin{tabular}{@{}p{0.32\textwidth}p{0.64\textwidth}@{}}
\toprule
\textbf{Decisión} & \textbf{Motivo} \\
\midrule
\textbf{Python 3, solo stdlib} (sin \cod{pip}) &
En las VMs no hay internet garantizado; \cod{socket}, \cod{email} y
\cod{http.server} bastan. Se arranca con \cod{py -3 servidor.py} en Windows. \\
\textbf{SMTP desde cero} (RFC 5321) &
Es lo didáctico y lo evaluado: HELO/EHLO, MAIL, RCPT, DATA y los códigos
2xx/4xx/5xx. Una librería (\cod{aiosmtpd}) ocultaría el protocolo. \\
\textbf{Webmail con \cod{http.server}} &
Cubre lo esencial de un <<Gmail>> (bandeja, leer, descargar adjuntos,
redactar) sin Flask ni dependencias. \\
\textbf{Sin autenticación} &
Para simplificar la complejidad de este ejercicio, como estamos en un entorno simulado hemos decidido \textbf{no} implementar un sistema de autenticación. \\
\textbf{Servidor en una VM Win10/11 dedicada} &
Puertos 25 (SMTP) y 8080 (webmail) en la red NAT de VirtualBox ya
configurada. \\
\textbf{Almacenamiento en ficheros} &
\cod{buzones/<usuario>/mensajes/*.eml} + \cod{adjuntos/} extraídos: fácil de
inspeccionar como evidencia. \\
\textbf{Pruebas con \cod{unittest}} (stdlib) &
Para verificar que el sistema funcionad e manera correcta.\\
\bottomrule
\end{tabular}
\caption{Decisiones técnicas del servidor de correo.}
\label{tab:decisiones-servidor}
\end{table}


% ---------------------------------------------------------------------
\section{Almacenamiento  de Logs y evidencias (cadena de custodia)}
Una vez hemos expuesto las decisiones de diseño de nuestro servidor vamos a detaller como funciona el registro de correos, ya que es la parte más importante del sistema para poder sacar conclusiones en un proceso de peritaje informático.\\

\begin{itemize}
  \item Cada mensaje se guarda como \textbf{\cod{.eml}} (mensaje MIME completo
        con todas las cabeceras, incluida \cod{Received} con la IP de origen,
        el HELO y la fecha).
  \item Los \textbf{adjuntos se extraen} a una carpeta por mensaje, con
        nombres saneados, de modo que un \cod{.pdf} se abre sin parsear el
        \cod{.eml}.
  \item El identificador de cola (\cod{cola\_id}) tiene resolución de segundo +
        UUID (\textit{Universally Unique Identifier}), el identificador único de cada correo.
  \item \cod{usuarios.json} guarda la lista de buzones (el apartado donde se almacenan los coreos de cada usuario.) Siguiendo el siguiente formato(\cod{{"usuarios": [...]}}).
\end{itemize}

Con esto conseguimos que la evidencia se recopile \textbf{sin ser modificada} y queda lista para la cadena de custodia.

\subsection{Ejemplo de registro (\textit{log}) de un envío de correo}
\begin{lstlisting}[style=terminal,breaklines=true]
	2026-11-05T10:00:00 | de: pepe@empresa.local | para: ana@empresa.local, luis@empresa.local | asunto: Examen final | adjuntos: Confidencial.pdf (123 KB), notas.txt (4 KB)
\end{lstlisting}

Como podemos ver cada mensaje entregado, ya sea por \textbf{SMTP o por el webmail}, añade \textbf{una única línea} a \cod{logs/envios.log} (en lenguaje mundano, se guarda en la carpeta logs en el fichero envios.). 
Entendamos ahora un poco qué información nos proporciona este registro:
\begin{itemize}
	\item Timestamp \href{https://www.iso.org/standard/40130.html}{ISO 8601} en hora local (segundos).
	\item \cod{para}: destinatarios en el orden aceptado (RCPT TO / lista del
	formulario).
	\item \cod{adjuntos}: nombres saneados con tamaño redondeado a KB (mínimo
	1 KB); si no hay, \cod{adjuntos: ninguno}.
	\item \textbf{Escaping}: si el asunto o un nombre contiene \cod{ | } (o
	cualquier \cod{|}), se sustituye por \textbrokenbar{} (U+00A6), de
	modo que cada envío es \textbf{exactamente una línea} jamás rota.
	\item El fichero \textbf{crece por anexión} (nunca se reescribe); la
	escritura va protegida con lock y \cod{try/except}: un fallo del log
	(disco lleno, permisos\ldots) queda en el logger
	\cod{correo.log_envios} pero \textbf{nunca tumba} el envío ni la
	sesión.
\end{itemize}

% ---------------------------------------------------------------------
\section{Por qué esto se adapta al escenario de la fuga por correo}

\begin{enumerate}
  \item Arrancar \cod{todo} en la VM Win10/11 del servidor de correo.
  \item Crear dos buzones: \cod{pepe@empresa.local} y
        \cod{ana@empresa.local}.
  \item En el Windows Server 2008: enviar \cod{Confidencial.pdf} adjunto con
        cualquier cliente SMTP (o telnet/PowerShell) a
        \cod{ana@empresa.local}.
  \item Evidencias que quedan en el servidor:
        \begin{itemize}
          \item \cod{servidor/buzones/ana/mensajes/<id>.eml} --- mensaje
                completo con la cabecera \cod{Received} (origen, HELO,
                fecha);
          \item \cod{servidor/buzones/ana/adjuntos/<id>/Confidencial.pdf} ---
                el adjunto extraído, listo para imagen y cadena de custodia;
          \item \cod{servidor/logs/envios.log} --- línea de auditoría del
                envío.
        \end{itemize}
  \item En el webmail, bandeja de \cod{ana} $\rightarrow$ abrir $\rightarrow$
        descargar adjunto.
  \item Archivar todo (imágenes de disco) \textbf{sin modificar} y elaborar la
        cadena de custodia.
\end{enumerate}

Esto cubre exactamente lo que pide el enunciado: sacar \cod{Confidencial.pdf}
del Windows Server 2008, enviarlo por correo entre dos miembros del grupo y
conservar la evidencia analizable.


% ---------------------------------------------------------------------

